% =============================================================================
% Chapter 8: Conclusion and Outlook
% FAPS: at most three pages. Active voice permitted for the own contribution.
% Subjunctive permitted in the outlook only.
% =============================================================================

\chapter[Conclusion and Outlook]{Conclusion and Outlook on Indicator-Based Operational Safety Assurance}
\label{ch:conclusion}

% =============================================================================
\section{Conclusion}
\label{sec:conclusion}

Operational safety assurance requires that the claims of a safety case be kept
under observation after release, and the instrument of that observation is the
safety performance indicator. UL~4600 \cite{ul4600} and PAS~1881
\cite{pas1881} require such monitoring and give examples of indicators;
ISO~21448 \cite{iso21448} requires that functional insufficiencies be monitored
in operation without saying which quantities carry that information; and the
assurance frameworks that consume indicator values take the indicators as given.
The step between an unresolved doubt in a safety case and a named signal with a
source and a threshold was the gap this thesis addressed.

This thesis makes three contributions.

The first is the Causal Chain Derivation Method. It takes a lagging outcome that
a Goal Structuring Notation claim denies and returns a leading indicator with a
signal, a source, a metric threshold and a documented path back to the claim, in
seven traceable steps: a lagging-indicator register, a proximate cause selected
by a prevention test, a decomposition through gate logic, the earliest observable
signal in the formulation the geometry admits, an observability and independence
check, the two thresholds, and attachment to the goal that the indicator actually
monitors.

The second is an evaluation of the resulting register on 94,816 simulated
unprotected left turns across three campaigns, in which the violations of six
leading indicators are set against the 540 oncoming-vehicle, 815 pedestrian and
101 following-vehicle collisions of the same runs.

The third is the monitoring of all nine indicators against their SPI thresholds
in a dashboard built on the same results, which demonstrates the form in which a
derived register is consumed.

Table~\ref{tab:rq_answers} states what the evaluation establishes for each
research question.

% ---------------------------------------------------------------------------
\begin{table}[htbp]
    \centering
    \caption{Answers to the research questions. The table states the answer
    established by the evaluation and the evidence on which it rests.}
    \label{tab:rq_answers}
    \small
    \begin{tabular}{@{}l>{\raggedright\arraybackslash}p{1.6cm}>{\raggedright\arraybackslash}p{8.4cm}@{}}
        \toprule
        {\bfseries Research question} & {\bfseries Answer} & {\bfseries Evidence} \\
        \midrule
        RQ1 Derivation &
        Yes &
        Seven traceable steps produce the register, and the checks reject
        candidates: PLAN\_DEV as an input of the controller's own error, the
        pedestrian perception branch as architectural, and the follower reaction
        delay as not observable to the ego vehicle \\
        RQ2 Timeliness &
        Partly &
        The leading SPIs remain measurable where collisions are rare, with 527 to
        9,841 violations against one collision in 75,509 runs; but only REAR\_TTC
        warns at least one second ahead in most collisions of its chain, 87 of
        101 at a median lead of 1.15\,s \\
        RQ3 Scalability &
        Yes &
        Indicators were derived and evaluated for all three conflict chains; the
        crossing chains require the time to the closest point of approach in
        place of the time to collision \\
        \bottomrule
    \end{tabular}
\end{table}

The answer to the second question is the one that qualifies the others. In the
two crossing conflicts the median lead time is 0.60\,s for GAP\_TTC and 0.50\,s
for PED\_CPA, and at least one second of warning is reached in 38 of 540 and two
of 815 collisions respectively. Measured against an in-run response, the derived
indicators of the crossing chains are not timely. Measured against an assurance
process that reads the violation rate of many runs, the picture differs by
indicator: the violation rate of PED\_CPA in Campaign~1 is 8.25 times its rate in
Campaigns~2 and~3, against 1.56 for GAP\_TTC. Which of the two readings applies
depends on what the indicator is monitored for, and the register records it.

Three statements bound every result reported here. The evaluation validates the
derivation method and the evaluation itself, not the safety of a driving
function. The metric threshold marks unusual behaviour rather than unacceptable
behaviour, and the SPI threshold is statistical for the same reason. And where
collisions are too rare to count, as in the 75,509 runs of the realistic-driving
campaigns, the leading SPIs are what remains to be monitored.

% =============================================================================
\section{Outlook}
\label{sec:outlook}

Several of the limitations in Section~\ref{sec:disc_limitations} point to work
that could be carried out on the existing data.

A sensitivity analysis under the unrestricted pedestrian collision definition
would bound how far the pedestrian results depend on the moving-ego rule of
Section~\ref{sec:exp_collision}, and would belong in an appendix beside the
results reported here. A per-run check of the 110 AEB\_PED runs would settle
whether the interval of about five to six seconds reflects a stop-and-go sequence
or a genuine warning, and would decide whether the reaction tier of the
pedestrian chain carries any timing information at all. Both could be done
without new simulation.

Beyond the present data, the SPI threshold could be derived from a risk
acceptance principle rather than from a statistical baseline. Setting T from
ALARP, GAMAB, MEM or a careful and competent human driver benchmark would require
a stated level of acceptable risk, and would have to be done with safety
engineers rather than from the data alone.

The register could be extended to the architectural tier. The indicators
evaluated here all measure kinematics or the vehicle's reaction to them, whereas
the doubts an automated driving system's safety case most often leaves open
concern the components that build its world model. Deriving an indicator from a
doubt about a perception component would require a signal observable at runtime
and informative about that component's state, and a confidence-based candidate
would require the model's confidence to be calibrated before a threshold on it
could carry meaning.

The exposure could be widened. A gap-accepting driving function with responsive
conflict partners would produce oncoming conflicts that arise from a functional
insufficiency rather than from a missing capability, and naturalistic trajectory
data [CITATION NEEDED: inD dataset] [CITATION NEEDED: SIND dataset] would place
the same indicators under a distribution of operating conditions that a designed
parameter grid cannot reproduce. Extreme value theory and human-driver baselines
for the metric threshold would become available with that data, and a sensitivity
sweep over the calibration percentile would show how much the reported operating
points depend on the tenth and ninety-fifth percentiles.

Finally, the method itself could be tested rather than only applied. A study in
which several engineers derive indicators independently from the same claim would
show how far the seven steps constrain the result, which the present work cannot
establish from a single application. A cross-check against a systems-theoretic
analysis would show whether causes that arise only from interacting functions,
and that a gate decomposition of component failures cannot reach, are being
missed.